\documentclass[11pt]{article}
% Built with xelatex: the appendix quotes prompts in Russian, set in DejaVu Serif.
\usepackage{fontspec}
\newfontfamily\cyrfont{DejaVu Serif}[Scale=MatchLowercase]
\newcommand{\ru}[1]{{\cyrfont #1}}
\usepackage[a4paper,margin=0.85in]{geometry}
\usepackage{hyperref}
\usepackage{enumitem}
\usepackage{parskip}
\usepackage{booktabs}
\usepackage{tabularx}

\hypersetup{colorlinks=true,urlcolor=blue,linkcolor=black}
\setlist{nosep,leftmargin=*}
\setlength{\parskip}{0.45em}

\title{\vspace{-3em}CS5013 Principles \& Practices\\[0.2em]
\large Exchange Student Guide \& Community Platform}
\author{Mikhail Novikov, GE26Z858 \and Abdirakhim Ismailov, GE26Z860}
\date{\vspace{-0.6em}9 October 2026}

\begin{document}
\sloppy
\maketitle
\vspace{-2.2em}

\noindent Repository: \url{https://github.com/rikire/Exchange-Student-Guide-Platform}. The code is
written with Claude Code. What follows is how we keep that honest: each principle names the file or
mechanism that enforces it and a dated case from the history where it mattered. The prompt journal
(\texttt{docs/ai/journal/}, 1087 entries since 3 September) holds every prompt quoted here.

\section{Prompt engineering}

\textbf{Context in layers, not one long prompt.} Standing rules live in files the assistant loads
itself: \texttt{CLAUDE.md} (six rules and where to look), \texttt{docs/ai/*.md} (one topic each:
prompting, testing, when to stop) and \texttt{.claude/rules/} (the short part that must be in view
while code is written). A prompt then carries only the task. A rule read once at the start of a
session loses to everything said later, so a \texttt{UserPromptSubmit} hook re-injects the core
rule with every prompt.

\textbf{A contract before any change.} For a request that changes files or state, the assistant
first reads what the repository can answer, then states the goal, the boundaries, acceptance criteria
phrased so they can become test names, open questions \emph{each with a suggested answer}, and steps
each with its check --- and waits for an explicit ``yes'' (\texttt{docs/ai/prompting.md}). ``Propose
and continue'' is not agreement; choosing a default and announcing it is not compliance. The
confirmed sentences become the acceptance criteria and the test names, so one sentence runs from
the prompt to the assertion. Appendix A.1 and A.3.

\textbf{Closed questions with a recommendation.} A question to the human offers options, the
consequence of each, and which one we would pick, so the answer can be one word (``A'', ``yes'').
Facts the repository can answer are looked up, never asked.

\textbf{Make the model verify, not recall.} Versions are checked against Maven Central or npm before
they are written down; a claim about a library is checked in its source. Saying ``not verified'' is
allowed; a plausible guess is not. When the human challenged a claim (``what is this ban on
JavaScript?'', A.2), the answer was a search of the requirements, not a defence of the claim.

\textbf{Reusable prompts as skills.} Recurring tasks are slash commands in
\texttt{.claude/skills/}: \texttt{/sharpen} writes the contract without acting, \texttt{/feature},
\texttt{/gaps}, \texttt{/dod}, \texttt{/trace-check}.

\section{Debugging}

We follow the order in \texttt{docs/ai/stop-and-ask.md}: reproduce, form one hypothesis that a
check can refute, test it, and fix with a test that was red first. Observable triggers stop the work:
a check that refutes a hypothesis, a second attempt with nothing new learned, a result that
contradicts a requirement. After a stop we write down what is confirmed, refuted and unknown, and
trace everything that rested on the wrong premise.

\textbf{A case: the moderator's approval answered 500 (5 October).} Seen only in a rehearsal of the
demo on the Docker stand. \emph{Evidence}: the application log, not the page. \emph{Hypothesis}:
the search index on the stand's volume predated a change to its fields. \emph{Check}: the database
held the approval while search did not find it, and the index builder skipped any index that held
documents. \emph{Fix}: rebuild the index at every start, an ADR amended with the human's decision,
and a test inverted to state it. The same rehearsal found three more faults no unit test could see:
volumes the earlier root image had written (\texttt{DEBT-026}), the missing error page, and a list
of tag suggestions covering the Submit button --- each fixed with the test that would have caught it.

\textbf{Read the library, then prove the reading.} Before adding an image-preview plugin we read its
bundle and found it loads images from \texttt{blob:} URLs, which our Content-Security-Policy blocks
(A.3). After widening the policy we restored the old one on purpose and watched the browser test go
red with the browser's own CSP message, so the test is known to catch it.

\textbf{Do not close what is not proven.} A tag-field browser test failed twice in full runs and
never alone. Tom Select's source showed a 300\,ms refresh delay that explains it; the fix removed the
delay, but no test could reproduce the failure first, so \texttt{DEBT-023} stays open with the
finding and the trigger for reopening it.

\section{Version control: code and prompts}

\textbf{The prompts are in Git.} Hooks write every prompt, an English rendering when it was not in
English, the outcome, the human's own edits between turns and the checks that passed, into
\texttt{docs/ai/journal/}; the \texttt{Stop} hook commits the entry. The author is resolved from the
Git identity. The original prompt is kept beside the translation: it is the artefact.

\textbf{Commits say what they deliver.} \texttt{<type>(<scope>): <subject> [FR-\ldots]}; a
\texttt{feat} or \texttt{fix} must name its feature and requirement, checked by the
\texttt{commit-msg} hook. Code carries \texttt{//trace:FR-XXX} anchors, so
\texttt{docs/traceability.md} is generated from code, tests and commits, and a requirement with no
test shows as a gap.

\textbf{Checks that cannot be skipped.} \texttt{pre-commit}: Spotless and fast tests;
\texttt{pre-push}: the full build plus ``generated documents current'' and ``a changed controller is
described in the route contract''; CI repeats it and runs the suite on PostgreSQL. The client refuses
\texttt{-{}-no-verify} and \texttt{-DskipTests} outright. On 2 October a push was refused three times
--- stale traceability, stale gap list, an undocumented controller change --- and each was fixed,
not bypassed. Unpushed work is rebased onto the partner's (\texttt{pull.ff=only}), and conflicts in
generated files are resolved by running the generator, never by hand.

\section{Testing, refactoring and documentation}

\textbf{Test first.} Interfaces, a red test, the minimal code, then refactor
(\texttt{docs/ai/workflow.md}); a test green before the change is reported as a guard, not as
evidence. Test names are sentences
(\texttt{a\_refused\_file\_is\_marked\_beside\_the\_file\_field\_and\_must\_be\_chosen\_again}).
Today: 635 application tests and 352 for the tools; 51 browser tests in Chromium with axe at four
widths; query-count tests that fail when a page's queries grow with its rows; mutation testing
(PIT, 130 of 132 mutants killed, the two survivors explained); property-based tests (jqwik, 1000
generated titles in English, Hindi and Tamil).

\textbf{Refactoring toward one place for each rule.} When the tag field needed to compare typed tags
with stored ones, the normalisation rule was moved out to \texttt{Tags.storedAs} rather than copied:
a copy of that rule had already been debt once (\texttt{DEBT-005}). Spring Modulith and ArchUnit
tests keep the thirteen slices from reaching into each other's internals.

\textbf{Documentation that cannot drift.} Changed behaviour updates its document in the same commit
(\texttt{docs/ai/docs-sync.md}); a \texttt{PostToolUse} hook names the document an edit puts out of
date. Four documents are generated and refused at push when stale: traceability, the gap list, the
route tables, slice ownership. Decisions are ADRs (22, each with the options weighed); anything
temporary is a \texttt{DEBT} entry with cause, consequence, fix and trigger, and a \texttt{TODO}
without one is refused.

\clearpage
\appendix
\section*{Appendix A. Sample prompts, from the journal}

Each prompt is quoted as typed (Russian), with the English rendering the journal keeps and an
abridged outcome.

\textbf{A.1 --- A short prompt, a contract before state changes} (\texttt{2026-10-02}, 14:50)\\
\emph{Prompt:} ``\ru{сделай пулл, проверь что сделано и что нужно сделать}'' --- \emph{Do a pull, check
what has been done and what needs to be done.}\\
\emph{Outcome:} fetched 125 commits; did not pull, since the branch had diverged and
\texttt{pull.ff=only} would refuse; proposed \texttt{pull -{}-rebase -{}-autostash} and waited.
\emph{Next prompt:} ``\ru{да, делай rebase}'' --- \emph{Yes, do the rebase.}

\textbf{A.2 --- Challenging the model's claim} (\texttt{2026-10-02})\\
\emph{Prompt:} ``\ru{погоди, а что за запрет на js? откуда он вообще и зачем}'' --- \emph{Wait, what is
this ban on JS? Where does it come from and why?}\\
\emph{Outcome:} checked the requirements and ADRs: nothing forbids JavaScript; the ``rule'' was one
descriptive sentence in a feature file. The assistant said it had overstated it and changed its
recommendation.

\textbf{A.3 --- A stop on a trigger} (\texttt{2026-10-02})\\
\emph{Prompt:} ``\ru{1. давай твой вариант. начинай}'' --- \emph{Go with your option. Start.}\\
\emph{Outcome:} reading the image-preview plugin's bundle against the CSP refuted the assumption that
it needed no policy change. The assistant stopped before writing the ADR and asked: widen the policy
by \texttt{img-src blob:} and \texttt{worker-src blob:}, or drop the preview. The human chose the
first.

\textbf{A.4 --- A contract for a multi-step change} (\texttt{2026-10-02})\\
The plan for the submission forms, as the assistant proposed it and the human confirmed: seven steps
(an ADR; the tag field; the drop zone; the editor; errors at the field; the draft; the edit form),
each with its check, one commit each, ``each step in turn, together''. Open questions came with a
suggested answer, e.g. \emph{exclude the library's transitive WebJars? --- suggested: yes, the
bundle contains them}.

\textbf{A.5 --- A debugging prompt} (\texttt{2026-10-05})\\
\emph{Prompt:} a screenshot of \emph{HTTP Status 500} and ``\ru{после аппрува}'' --- \emph{after the
approval.}\\
\emph{Outcome:} read the stand's log; found the approval saved and the index write refused by a
mapping from an older version; offered three ways forward with a recommendation; the human answered
``\ru{1A, 2 да, 3 да, 4 да}'' --- \emph{1A, 2 yes, 3 yes, 4 yes}.

\end{document}
